% Generated draft: requires mathematical and source review.
\documentclass[11pt]{article}
\usepackage[margin=25mm]{geometry}
\usepackage{fontspec}
\setmainfont{Latin Modern Roman}
\usepackage{amsmath,amssymb,mathtools,unicode-math}
\setmathfont{Latin Modern Math}
\usepackage{enumitem,needspace,xurl,hyperref}
\setlength{\parindent}{0pt}
\setlength{\parskip}{5pt}
\setlength{\emergencystretch}{3em}
\newcommand{\N}{\mathbb{N}}
\newcommand{\Z}{\mathbb{Z}}
\newcommand{\Q}{\mathbb{Q}}
\newcommand{\R}{\mathbb{R}}
\newcommand{\C}{\mathbb{C}}
\newcommand{\F}{\mathbb{F}}
\newcommand{\A}{\mathbb{A}}
\newcommand{\PP}{\mathbb{P}}
\newcommand{\E}{\mathbb{E}}
\newcommand{\eps}{\varepsilon}
\newcommand{\dd}{\,\mathrm{d}}
\begin{document}
\section*{UnsolvedMath ID: 30000498}
\textbf{Source problem code:} catalogue entry
\subsection{Definitions and mathematical statement}
\textbf{Weak Greenberg conjecture, curated form.} For every odd prime \(p\) and every CM field \(K\) containing \(\mu_p\), is\[\lambda^-(K)>\lambda^+(K)\quad\text{unless }\lambda^-(K)=\lambda^+(K)=0?\] A CM field is a totally imaginary quadratic extension of a totally real field \(K^+\); its complex conjugation is denoted \(c\). The condition \(\mu_p\subset K\) means that all \(p\)th roots of unity belong to \(K\). Neither the odd-prime hypothesis nor this root-of-unity hypothesis is removed.\par Under the standard cyclotomic Iwasawa convention, let \(K_\infty/K\) be the cyclotomic \(\Z_p\)-extension, with layers \([K_n:K]=p^n\). Set \(A_n=\operatorname{Cl}(K_n)[p^\infty]\) and \(X=\varprojlim_n A_n\), with norm transition maps. Complex conjugation acts on these groups; since \(p\ne2\), the idempotents \((1+c)/2\) and \((1-c)/2\) define the plus and minus parts \(X^+,X^-\). The invariants \(\lambda^+,\lambda^-\) are the Iwasawa \(\lambda\)-invariants of these torsion modules over \(\Z_p[[\operatorname{Gal}(K_\infty/K)]]\), equivalently the total degrees of their distinguished characteristic-polynomial factors with multiplicity. This is an explicit conventional interpretation: the excerpt does not itself define the tower, module, or invariant normalization. No vanishing of the corresponding \(\mu\)-invariants is assumed merely in order to ask the question.\par The original prose instead states \(\lambda^+=\lambda^-\) if and only if both vanish. The curated strict-inequality form also rules out \(\lambda^-<\lambda^+\). Identifying these formulations therefore requires the relevant comparison inequality under the same hypotheses, which is not supplied as a theorem in this record. Both formulations are preserved rather than declared logically equivalent without that input. The weak conjecture does not by itself assert \(\lambda^+=0\); strict inequality with a positive plus invariant remains allowed.\par The supplied triage checked 2026-08-21 labels the problem partially solved. It reports substantial special hypotheses, described as abelian plus fields with total splitting and suitable circular-unit cohomology, while saying the general assertion remains open. The exact field, splitting, and cohomological conditions are not reproduced, so no unconditional theorem for all abelian CM fields is inferred. This dated assessment requires independent verification.
\subsection{Short English statement}
For the cyclotomic plus and minus Iwasawa invariants of a CM field containing \(\mu_p\), the curated claim requires strict minus dominance except when both vanish. The original equality-only statement needs a comparison inequality to be equivalent. Neither formulation is silently strengthened to vanishing of the plus invariant.
\subsection{Sources}
\begin{itemize}\item Local UnsolvedMath registry, permanent ID 30000498, code catalogue entry; \emph{Weak Greenberg Conjecture for CM Fields}, attributed to Greenberg in the original statement; Oberwolfach Reports Open Problems.\item \emph{Pro-\(p\) Extensions of Global Fields and pro-\(p\) Groups (2007)}, Oberwolfach Reports, DOI 10.4171/owr/2006/25, pp. 1463--1536. \url{https://doi.org/10.4171/owr/2006/25}. Proposed year and year tag: 2006; no separate proposer field supplied. The source's escaped title delimiters are rendered as mathematics.\item Supplied literature label \emph{original OWR report}: \url{https://doi.org/10.4171/owr/2006/25}.\item Supplied literature label \emph{Nguyen Quang Do primary paper}: \url{https://doi.org/10.1142/S1793042106000395}.\item Supplied literature label \emph{later Greenberg-conjecture framework}: \url{https://arxiv.org/abs/1612.00718}.\item Literature triage checked 2026-08-21, status \texttt{partially\_solved}; record updated 2026-09-16. The \texttt{corrected\_verified} field reports restoring odd-prime and strict-inequality hypotheses from the official report, not independently checked here.\end{itemize}
\end{document}
